% 土星（综述）
% license CCBYSA3
% type Wiki

本文根据 CC-BY-SA 协议转载翻译自维基百科\href{https://en.wikipedia.org/wiki/Saturn}{相关文章}。

\begin{figure}[ht]
\centering
\includegraphics[width=8cm]{./figures/13ddb48d3e7de073.png}
\caption{土星及其显著的行星环，由卡西尼轨道器拍摄\(^\text{[a]}\)} \label{fig_Saturn_1}
\end{figure}
土星是距离太阳第六颗行星，也是太阳系中仅次于木星的第二大行星。它是一颗气态巨行星，其平均半径约为地球的 9 倍。它的平均密度仅为地球的八分之一，但质量却超过地球的 95 倍。尽管土星几乎与木星一样大，但其质量还不到木星的三分之一。土星以 9.59 AU（14.34 亿公里）的距离绕太阳运行，轨道周期为 29.45 年。

土星内部被认为由一颗岩石核心组成，其外包裹着一层深厚的金属氢层，其上为液态氢和液态氦的中间层，最外层为气体外壳。土星呈现淡黄色调，是由于其上层大气中的氨晶体所致。金属氢层中的电流被认为会产生土星的行星磁场，该磁场强度弱于地球，但由于土星体积更大，其磁矩是地球的 580 倍。土星磁场的强度约为木星的二十分之一\(^\text{[27]}\)。土星外层大气整体较为平淡、缺乏对比度，但亦可能出现长期存在的结构。土星的风速可达每小时 1,800 公里（1,100 英里）。

这颗行星拥有明亮且庞大的光环系统，主要由冰粒构成，另含少量岩石碎屑与尘埃。至少有 274 颗卫星绕火星运行，其中 63 颗已有正式命名；这一数字不包含光环中的数百颗小卫星体。土卫六（Titan）是土星最大的卫星，也是太阳系中第二大的卫星，其体积大于（但质量小于）水星，并且是太阳系中唯一拥有稠密大气层的卫星\(^\text{[28]}\)。
\subsection{名称与符号}
土星以罗马财富与农业之神命名，该神是木星（Jupiter）之父。其天文符号（♄）可追溯至希腊的《奥克绪林库斯纸草文稿》（Oxyrhynchus Papyri），在其中该符号可见为希腊字母 κ–ρ（kappa–rho）的连写并带有一条水平短线，作为 Κρονος（Cronus，即土星的希腊名称）的缩写\(^\text{[29]}\)。该符号后演化出类似小写希腊字母 eta 的形状，并在 16 世纪顶部添加十字，以使这一原本的异教符号基督化。

罗马人以土星之名将一周的第七天命名为 Saturday，即 Sāturni diēs，“土星之日”\(^\text{[30]}\)。
\subsection{物理特征}
\begin{figure}[ht]
\centering
\includegraphics[width=8cm]{./figures/91b99dbcbbce021d.png}
\caption{土星与地球及地球月球大小的比较} \label{fig_Saturn_2}
\end{figure}
土星是一颗气态巨行星，主要由氢和氦组成。它缺乏明确的固体表面，但行星内部很可能存在一颗固体核心\(^\text{[31]}\)。由于自转的影响，土星呈现扁球体形状——两极被压扁、赤道区域鼓起。其赤道半径比极半径大超过 10\%：分别为 60,268 公里与 54,364 公里（37,449 英里与 33,780 英里）\(^\text{[6]}\)。

太阳系中的其他巨行星——木星、天王星与海王星——其扁率均低于土星。赤道凸起与自转速率的共同作用导致土星赤道处的有效表面重力仅为 8.96 m/s²，相当于其两极重力的 74\%，且低于地球表面重力。然而，赤道处的逃逸速度接近 36 km/s，远高于地球的逃逸速度\(^\text{[32]}\)。

土星是太阳系中唯一密度低于水的行星——约低 30\%\(^\text{[33]}\)。尽管其核心的密度远高于水，但由于其厚重大气层的存在，整颗行星的平均密度仅为 0.69 g/cm³。木星的质量是地球的 318 倍\(^\text{[34]}\)，而土星的质量是地球的 95 倍\(^\text{[6][35][36][37][38][39]}\)。木星与土星合计占据了太阳系全部行星质量的 92\%\(^\text{[40]}\)。
\subsubsection{内部结构}
\begin{figure}[ht]
\centering
\includegraphics[width=8cm]{./figures/0e870b2d45e0c034.png}
\caption{按比例绘制的土星示意图} \label{fig_Saturn_3}
\end{figure}
尽管土星主要由氢和氦组成，但其大部分质量并非处于气态，因为当密度超过 0.01 g/cm³ 时，氢会成为非理想液体，而这一密度在包含土星 99.9\% 质量的半径范围内即可达到。土星内部的温度、压力与密度在向核心方向不断上升，使得在更深层中氢呈金属态\(^\text{[40]}\)。

标准的行星模型表明，土星内部与木星相似，均包含一颗小型岩石核心，外围包裹着氢和氦，并含有少量各种挥发物\(^\text{[41]}\)。对土星形变的分析显示，其中心物质的集中程度显著高于木星，因此在其中心附近含有更多密度高于氢的物质。土星中心区域按质量计算约有 50\% 为氢，而木星约为 67\%\(^\text{[42]}\)。

该核心在组成上与地球类似，但密度更大。结合土星的引力矩测量结果与其内部物理模型，科学家得以对土星核心的质量设定约束。2004 年的估计认为其核心质量必须为地球质量的 9–22 倍\(^\text{[43][44]}\)，相当于直径约 25,000 公里（16,000 英里）\(^\text{[45]}\)。对土星光环的测量则显示其核心可能更加弥散，质量约等于 17 个地球，半径约为土星整体半径的 60\%\(^\text{[46]}\)。其外层包裹着更厚的液态金属氢层，然后是被氦饱和的液态分子氢层，随高度升高逐渐过渡为气体。最外层约有 1,000 公里（620 英里）厚，由气体组成\(^\text{[47][48][49]}\)。

土星内部非常炽热，其核心温度可达 11,700 °C（21,100 °F），并向外辐射的能量是其自太阳接收能量的 2.5 倍。木星的热能由缓慢的引力压缩（开尔文–亥姆霍兹机制）产生；但这一机制单独可能不足以解释土星的热量来源，因为土星质量更小。另一种或额外的机制可能是土星深层内部的氦滴“降雨”过程。当氦滴向密度较低的氢层下降时，摩擦会释放能量，并使土星外层的氦含量减少\(^\text{[50][51]}\)。这些下降的氦滴可能最终形成包围核心的氦层\(^\text{[41]}\)。在土星内部，如同在木星\(^\text{[52]}\)、以及冰巨行星天王星和海王星\(^\text{[53]}\)中一样，可能会出现钻石雨。
\subsubsection{大气层}  
土星外层大气按体积分数包含 96.3\% 的分子氢和 3.25\% 的氦。与太阳中该元素的丰度相比，土星大气中的氦比例显著偏低\(^\text{[41]}\)。比氦更重的元素（金属丰度）的精确含量尚不明确，但通常认为其比例与太阳系形成时期的原始丰度一致。估计这些较重元素的总质量为地球质量的 19–31 倍，其中相当一部分位于土星的核心区域\(^\text{[54]}\)。

在土星大气中已探测到微量的氨、乙炔、乙烷、丙烷、磷化氢和甲烷\(^\text{[55][56][57]}\)。上层云由氨晶体组成，而下层云似乎由硫氢化铵（NH\(_4\)SH）或水构成\(^\text{[58]}\)。太阳的紫外辐射会在上层大气中引发甲烷的光解，导致一系列烃类化学反应，其生成物再通过湍流与扩散向下输送。此光化学循环受土星年度季节周期的调制\(^\text{[57]}\)。卡西尼号曾观测到一系列出现在北纬的云层结构，被昵称为“珍珠串”。这些结构是位于更深云层中的云隙\(^\text{[59]}\)。

\textbf{云层}
\begin{figure}[ht]
\centering
\includegraphics[width=8cm]{./figures/5b480d0d41a65c33.png}
\caption{2011 年，一场全球性的风暴环绕整个行星。该风暴沿着行星一周移动，使得风暴的头部（明亮区域）绕行并经过其尾部。} \label{fig_Saturn_4}
\end{figure}
土星的大气呈现出类似木星的带状结构，但土星的云带要更为暗淡，并且在赤道附近宽度更大。用于描述这些云带的命名法与木星相同。直到 20 世纪 80 年代“旅行者”号飞掠任务之后，人们才观察到土星更细致的云图案。自那以后，地基望远镜观测技术已有显著提升，使得对土星的定期观测成为可能。(^\text{[60]})

云层的组成随深度和压力的增加而变化。在上层云层中，温度位于 100–160 K 范围内，压力约为 0.5–2 bar，云层由氨冰组成。水冰云层开始于约 2.5 bar 的压力位置，并向下延伸至 9.5 bar，此处的温度范围为 185–270 K。在这一层中混杂有一条硫氢化铵冰带，处于 3–6 bar 的压力范围内，温度为 190–235 K。更低的位置，压力在 10–20 bar、温度 270–330 K 的区域，则含有由水滴和溶于水中的氨构成的混合层。(^\text{[61]})

通常较为平淡的土星大气偶尔也会呈现出类似木星的长期存在的椭圆云以及其他特征。1990 年，“哈勃”太空望远镜拍摄到赤道附近一片巨大的白色云团，而在“旅行者”号飞掠期间这一特征并不存在；1994 年又观测到另一场较小的风暴。1990 年的风暴是“大白斑”的一个典型例子，这是一种短暂现象，约每土星年（约合地球 30 年）在北半球夏至前后出现一次。(^\text{[62]})

早期的“大白斑”曾分别在 1876 年、1903 年、1933 年和 1960 年被观测到，其中 1933 年的风暴观测最为详细。(^\text{[63]}) 最近一次大型风暴出现在 2010 年。2015 年，研究者利用甚大阵列（VLA）望远镜研究了土星大气，报告称他们发现了“所有中纬度大型风暴的长存痕迹、存在数百年历史的赤道风暴混合体，以及可能在北纬 70° 处存在一场未被记录的更古老的风暴”。(^\text{[64]})

土星的风速是太阳系行星中第二快的，仅次于海王星。“旅行者”号的数据表明其东风峰值风速可达 500 m/s（1,800 km/h）。(^\text{[65]}) 在“卡西尼”号于 2007 年拍摄的图像中，土星北半球呈现出类似天王星的亮蓝色，这种颜色很可能由瑞利散射造成。(^\text{[66]}) 热成像显示土星南极存在一个温暖的极涡，这是太阳系中唯一已知的此类现象。(^\text{[67]}) 虽然土星的典型温度约为 −185 °C，但极涡中的温度经常可达 −122 °C，被认为是土星上最温暖的区域。(^\text{[67]})

\textbf{六边形云层结构}
\begin{figure}[ht]
\centering
\includegraphics[width=8cm]{./figures/77d74f1601865d59.png}
\caption{土星南北两极的红外图像} \label{fig_Saturn_5}
\end{figure}
在约 78°N 纬度处，围绕北极涡旋的大气中存在一个持久的六边形波形结构，最早在旅行者号的图像中被发现\(^\text{[68][69][70]}\)。六边形的每一条边约为 14{,}500\,\text{km}（9{,}000\,\text{mi}），比地球的直径还要长\(^\text{[71]}\)。整个结构以 10\,\text{h}\,39\,\text{m}\,24\,\text{s} 的周期旋转（与该行星的无线电辐射周期一致），这一周期被认为等同于土星内部的自转周期\(^\text{[72]}\)。与土星大气中其他可见云层不同，这一六边形结构在经度方向上并不发生漂移\(^\text{[73]}\)。这一图案的成因仍是高度推测的话题。多数科学家认为它是大气中的驻波结构。通过差速旋转流体的实验，可以在实验室中复制多边形形状\(^\text{[74][75]}\)。

哈勃太空望远镜对南极区域的成像表明，该区域存在一条急流，但并未发现强极涡或任何六边形驻波\(^\text{[76]}\)。NASA 于 2006 年 11 月报道，卡西尼号观测到一个“类飓风”风暴，其中心固定在南极，并具有一个清晰的风眼墙\(^\text{[77][78]}\)。在地球之外的任何其他行星上此前均未观测到风眼墙。例如，加利略号拍摄的木星大红斑图像中并未显示风眼墙\(^\text{[79]}\)。

南极风暴可能已经存在了数十亿年\(^\text{[80]}\)。该极涡的规模可与地球相当，风速可达 550\,\text{km/h}\(^\text{[80]}\)。
\subsubsection{磁层}
\begin{figure}[ht]
\centering
\includegraphics[width=8cm]{./figures/d9afeaaa48604fc3.png}
\caption{土星北极的极光} \label{fig_Saturn_6}
\end{figure}
土星具有一个固有的磁场，其形状简单且对称——即一个磁偶极子。其在赤道处的强度为 0.2 高斯（20~μT），大约是木星磁场强度的二十分之一，并且比地球的磁场稍弱\(^\text{[27]}\)。因此，土星的磁层远小于木星的磁层\(^\text{[81]}\)。

当旅行者 2 号进入磁层时，太阳风压较高，磁层仅延伸到 19 个土星半径处，即 110 万千米（684,000 英里）\(^\text{[82]}\)，但在几个小时内磁层扩张，并保持这种状态约三天\(^\text{[83]}\)。磁场很可能与木星的磁场生成方式相似——由液态金属氢层中的电流（称为金属氢发电机）所产生\(^\text{[81]}\)。该磁层在偏转来自太阳的太阳风粒子方面十分有效。土星的卫星泰坦在土星磁层的外部区域轨道运行，并通过其外层大气中的电离粒子向磁层补充等离子体\(^\text{[27]}\)。与地球类似，土星的磁层也会产生极光\(^\text{[84]}\)。
\subsection{轨道与自转}
\begin{figure}[ht]
\centering
\includegraphics[width=8cm]{./figures/4ee15093919e1056.png}
\caption{土星及太阳系外行星绕太阳运行的动画} \label{fig_Saturn_7}
\end{figure}
\begin{figure}[ht]
\centering
\includegraphics[width=8cm]{./figures/23c0c3dad6383630.png}
\caption{土星在一次轨道运行周期中（2001–2029 年），从地球（在冲日位置）观测到的模拟外观} \label{fig_Saturn_8}
\end{figure}
土星与太阳之间的平均距离超过 14 亿千米（9 AU）。其平均轨道速度为 9.68 km/s，\(^\text{[6]}\) 土星需要 10,759 个地球日（约 29\(\tfrac{1}{2}\) 年）\(^\text{[85]}\) 才能完成一次绕太阳的公转。\(^\text{[6]}\) 因此，它与木星形成近似的 5:2 平均运动共振。\(^\text{[86]}\) 土星的椭圆轨道相对于地球轨道平面倾斜 2.48°。\(^\text{[6]}\) 其近日点与远日点距离平均分别为 9.195 与 9.957 AU。\(^\text{[6][87]}\) 土星上可见的大气特征会随纬度不同以不同速率自转，因此与木星类似，不同区域被分配了不同的自转周期。

天文学家使用三套不同的系统来描述土星的自转速率。系统 I 的自转周期为 10h 14m 00s（844.3°/d），包括赤道带、南赤道带与北赤道带。极区自转速率被认为与系统 I 类似。除南北极区以外的其他土星纬度称为系统 II，其自转周期被指定为 10h 38m 25.4s（810.76°/d）。系统 III 则指土星内部的自转速率。根据旅行者 1 号与旅行者 2 号探测到的土星无线电辐射，\(^\text{[88]}\) 系统 III 的自转周期为 10h 39m 22.4s（810.8°/d）。系统 III 已在很大程度上取代了系统 II。\(^\text{[89]}\)

土星内部自转周期的精确数值仍难以确定。卡西尼号于 2004 年接近土星时发现土星无线电自转周期明显增加，约为 10h 45m 45s ± 36s。\(^\text{[90][91]}\) 结合卡西尼号、旅行者号与先驱者号的多项测量结果，对土星整体自转速率的估计为 10h 32m 35s。\(^\text{[92]}\) 对土星 C 环的研究得出的自转周期为 10h 33m 38s \(^{+1m 52s}_{-1m 19s}\)。\(^\text{[17][18]}\)

2007 年 3 月，人们发现土星无线电辐射的变化与其自转速率不符。这种差异可能由土卫二（Enceladus）上的喷泉活动引起。喷泉活动向土星轨道释放的水蒸气被电离后，会对土星磁场施加拖曳，使得磁场自转相对于行星本体自转略微减慢。\(^\text{[93][94][95]}\)

土星目前仅已知一颗特洛伊小行星，即 2019 UO\(_{14}\)，其特洛伊构型于 2024 年 9 月公布，围绕太阳运行于土星轨道前方 60° 的稳定 L\(_4\) 拉格朗日点。\(^\text{[96]}\) 这一发现使得水星成为唯一尚未发现特洛伊天体的行星。轨道共振机制，包括岁差共振，被认为是已知土星特洛伊天体数量极少的原因之一。\(^\text{[97]}\)
\subsection{天然卫星}
\begin{figure}[ht]
\centering
\includegraphics[width=10cm]{./figures/72eb16859e562019.png}
\caption{艺术家笔下的土星、其光环以及主要的冰质卫星——从米玛斯到瑞亚} \label{fig_Saturn_9}
\end{figure}
土星已知有 274 颗卫星，\(^\text{[98][99][100][101]}\) 其中 63 颗拥有正式名称。\(^\text{[12][11]}\) 有证据显示，在土星光环中存在直径 40–500 米的数十到数百个微小卫星体，\(^\text{[102]}\) 但它们不被视为真正的卫星。最大卫星——泰坦（Titan）——占据土星环绕系统（包括光环）总质量的 90\% 以上。\(^\text{[103]}\) 土星第二大卫星——瑞亚（Rhea）——可能拥有其自身的稀薄光环系统，\(^\text{[104]}\) 并具有极其稀薄的大气层。\(^\text{[105][106][107]}\)

其他许多卫星都很小：其中 131 颗的直径小于 50 公里。\(^\text{[108]}\) 传统上，土星的大多数卫星都以希腊神话中的泰坦命名。泰坦是太阳系中唯一拥有主要大气层的卫星，\(^\text{[109][110]}\) 其中存在复杂的有机化学反应。它也是唯一拥有烃类湖泊的卫星。\(^\text{[111][112]}\)

2013 年 6 月 6 日，IAA-CSIC 的科学家报告在泰坦上层大气中探测到多环芳香烃（PAHs），这可能是生命的前体物质。\(^\text{[113]}\) 2014 年 6 月 23 日，美国宇航局（NASA）宣布有力证据表明：泰坦大气层中的氮并非来自形成土星的原始物质，而是来自与彗星有关的奥尔特云物质。\(^\text{[114]}\)

土卫二（Enceladus）在化学成分上似乎与彗星类似，\(^\text{[115]}\) 因而长期被视为可能存在微生物生命的栖息地。\(^\text{[116][117][118][119]}\) 支持这一可能性的证据包括：其富含盐分的颗粒具有“类海洋”的化学组成，表明土卫二喷出的多数冰来自盐水液体蒸发。\(^\text{[120][121][122]}\) 卡西尼号在 2015 年穿越土卫二羽状喷流时，探测到维持产甲烷生命形式所需的主要成分。\(^\text{[123]}\)

2014 年 4 月，美国宇航局科学家报告称：土星 A 环中可能正在形成一颗新的卫星，此前卡西尼号在 2013 年 4 月 15 日的图像中拍摄到了该结构。\(^\text{[124]}\)
\subsection{行星环}
\begin{figure}[ht]
\centering
\includegraphics[width=8cm]{./figures/8c0274dc32d85fc0.png}
\caption{v土星的光环（此处为卡西尼号于 2004 年 10 月拍摄）是太阳系中最庞大、最显眼的光环系统。\(^\text{[48]}\)} \label{fig_Saturn_10}
\end{figure}
土星最为人所知的，或许就是其行星光环系统，使其在视觉上独具特色。\(^\text{[48]}\) 光环自土星赤道向外延伸，从 6,630 千米至 120,700 千米（4,120–75,000 英里），平均厚度约为 20 米（66 英尺）。光环主要由水冰组成，并含有微量的托林（tholin）杂质以及约 7\% 的无定形碳颗粒作为表面覆盖。\(^\text{[125]}\) 构成光环的颗粒大小从微小尘埃到直径 10 米不等。\(^\text{[126]}\) 虽然其他气体巨行星也拥有光环系统，但土星的光环最大且最为醒目。

关于光环的年龄存在争论。一种观点认为光环非常古老，与土星一起由原始星云物质形成（约 46 亿年前），\(^\text{[127]}\) 或形成于晚期重轰炸（LHB）之后不久（约 41–38 亿年前）。\(^\text{[128][129]}\) 另一种观点认为光环年轻得多，大约形成于 1 亿年前。\(^\text{[130][131][132]}\) 支持后者理论的麻省理工学院研究团队提出，光环是土星一颗被摧毁的卫星——名为“Chrysalis”——的残余物。\(^\text{[133]}\)

在主光环之外，在距离行星 1,200 万千米（750 万英里）处存在稀疏的菲比（Phoebe）光环。该光环相对于其他光环倾斜 27°，并像菲比一样以逆行方式运行。\(^\text{[134]}\)

土星的一些卫星，如潘多拉（Pandora）和普罗米修斯（Prometheus），作为“牧羊卫星”帮助限制光环、阻止其向外扩散。\(^\text{[135]}\) 潘（Pan）和阿特拉斯（Atlas）会在土星光环中激发弱的线性密度波，这些密度波已用于更可靠地计算它们的质量。\(^\text{[136]}\)
\begin{figure}[ht]
\centering
\includegraphics[width=14.25cm]{./figures/bf65e1a518a4adfe.png}
\caption{卡西尼号窄角相机于 2007 年 5 月 9 日拍摄的土星 D、C、B、A 与 F 环（从左至右）未受光照一侧的自然色拼接图像。图中距离均以土星中心为基准。} \label{fig_Saturn_11}
\end{figure}
\subsection{观测与探测的历史}
对土星的观测与探测可分为三个阶段：（1）前现代时期肉眼观测；（2）自 17 世纪起从地球进行的望远镜观测；（3）由空间探测器进行的访问，包括飞越与轨道探测。在 21 世纪，对土星的望远镜观测仍在继续（包括地球上的望远镜以及哈勃太空望远镜等地球轨道天文台），并由卡西尼号在土星轨道进行观测直至其于 2017 年退役。
\subsubsection{前望远镜时代的观测}  
参见：Saturn（神话）以及 占星术中的行星 § Saturn  
土星自史前时代以来就为人所知，\(^\text{[137]}\) 在早期的文字记录中，它是多种神话中的重要角色。巴比伦天文学家系统地观测并记录了土星的运动。\(^\text{[138]}\) 在古希腊语中，该行星被称为 Φαίνων Phainon，\(^\text{[139]}\) 而在罗马时代，它被称为“土星之星”或“太阳（即 Helios）之星”。\(^\text{[140][141]}\) 在古罗马神话中，Phainon 这颗行星被视为农业之神的神圣星体，现代名称即源于此。\(^\text{[142]}\) 罗马人认为土星神（Saturnus）与希腊神克罗诺斯（Cronus）相对应。在现代希腊语中，这颗行星仍保留克罗诺斯之名——Κρόνος：Kronos。\(^\text{[143]}\)

希腊科学家托勒密根据其在土星处于冲日位置时所做的观测，推算了土星的轨道。\(^\text{[144]}\) 在印度占星术中，有九大占星天体，被称为 Navagrahas。土星被称为 “Shani”，并根据人生中的善恶行为对每个人作出评判。\(^\text{[142][144]}\) 古代中国与日本文化将土星称为“土星”，源自用于分类自然元素的五行体系。\(^\text{[145][146][147]}\)

在希伯来语中，土星被称为 Shabbathai。\(^\text{[148]}\) 其守护天使为 Cassiel。其智慧体或善灵为 ‘Agȋȇl（希伯来语：אגיאל，拉丁化：ʿAgyal），\(^\text{[149]}\) 而其阴暗灵（恶魔）为 Zȃzȇl（希伯来语：זאזל，拉丁化：Zazl）。\(^\text{[149][150][151]}\) Zazel 被描述为在所罗门魔法中被召唤的伟大天使，“在爱情召唤中十分有效”。\(^\text{[152][153]}\) 在奥斯曼土耳其语、乌尔都语与马来语中，Zazel 的名字为 ‘Zuhal’，源自阿拉伯语（阿拉伯语：زحل，拉丁化：Zuhal）。\(^\text{[150]}\)
\subsubsection{望远镜时代的太空飞行前观测}
\begin{figure}[ht]
\centering
\includegraphics[width=8cm]{./figures/de1ff6ea3f6b74db.png}
\caption{罗伯特·胡克在 1666 年的这幅土星绘图中注意到球体与光环彼此投下的阴影（a 与 b）。} \label{fig_Saturn_12}
\end{figure}
\begin{figure}[ht]
\centering
\includegraphics[width=8cm]{./figures/b8d92d1f5c26ae9a.png}
\caption{伽利略·伽利莱在 1610 年观测到了土星的光环，但当时无法判断它们的本质。} \label{fig_Saturn_13}
\end{figure}
土星的光环至少需要一架口径 15 毫米的望远镜才能分辨，\(^\text{[154]}\) 在此之前它们并不为人所知，直到克里斯蒂安·惠更斯于 1655 年首次看到，并于 1659 年发表观测结果。伽利略在 1610 年用他原始的望远镜观测土星时，\(^\text{[155][156]}\) 错误地将土星看起来不像圆形的现象解释为其两侧各有一颗卫星。\(^\text{[157][158]}\)

当惠更斯使用更高的望远镜放大倍率时，这种看法被推翻，人类首次真正看清了土星光环。惠更斯还发现了土星的卫星——泰坦。后来，焦万尼·多梅尼科·卡西尼又发现了另外四颗卫星：土卫八（Iapetus）、土卫五（Rhea）、土卫三（Tethys）和土卫四（Dione）。1675 年，卡西尼发现了如今称为“卡西尼缝隙”的光环间隙。\(^\text{[159]}\)

此后直到 1789 年，没有进一步的重要发现，直到威廉·赫歇尔发现了另外两颗卫星：米玛斯（Mimas）和恩克拉多斯（Enceladus）。1848 年，一支英国团队发现了不规则形状的卫星海帕里翁（Hyperion），该卫星与泰坦存在轨道共振关系。\(^\text{[160]}\)

1899 年，威廉·亨利·皮克林发现了菲比（Phoebe），这是一颗高度不规则的卫星，与土星的大型卫星不同，它不是同步自转。\(^\text{[160]}\) 菲比是首颗被发现的此类卫星，它以逆行轨道绕土星运行，需要一年多才能完成一次环绕。20 世纪初，对泰坦的研究在 1944 年确认了它拥有厚重大气——这是太阳系众卫星中独一无二的特征。\(^\text{[161]}\)
\subsubsection{太空飞行任务}  
\textbf{先驱者 11 号飞掠}
\begin{figure}[ht]
\centering
\includegraphics[width=8cm]{./figures/b0f551c206af9e74.png}
\caption{先驱者 11 号拍摄的土星图像} \label{fig_Saturn_14}
\end{figure}
















